%%%
%%% Radical "⼲"
%%%
\section*{Radical 51: ``⼲''}\addcontentsline{toc}{section}{Radical 51: ⼲}\addcontentsline{loh}{figure}{\#\#\#\# 51: ⼲}

%%%%%%%%%% 干 %%%%%%%%%%
\subsection*{干}\addcontentsline{loh}{figure}{干}

\begin{Entry}{干}{3}{⼲}[Kangxi 51]
  \begin{Phonetics}{干}{gan1}
    \definition*{s.}{Sobrenome: Gan}
    \definition{adj.}{seco | vazio; oco; seco | sem substância; vazio | de parentesco nominal; (parentes) não ligados por laços sanguíneos | sem água; (água) esgotada; completamente vazia | assumido como parente nominal; relação familiar reconhecida por adoção | rude; grosseiro; mal"-educado; descreve alguém que fala de forma muito direta e rude (sem delicadeza).}
    \definition{adv.}{em vão; fútil; sem propósito; para nada; sem resultado | apenas; sem nada mais | inutilmente; sem uso, sem aproveitamento | superficialmente; significa que não há conteúdo, apenas forma}
    \definition{s.}{Arcaico: escudo | margem; ribeira; margem das águas | alimentos desidratados | abreviação para os dez troncos celestiais}
    \definition{v.}{ofender; afrontar | ter a ver com; estar relacionado com; estar implicado em; interferir com | (antiquado) buscar (cargo público, remuneração, etc.) | Dialeto: deixar alguém de fora; tratar alguém com indiferença; desprezar | assediar; perturbar; criar confusão; causar estragos; bagunçar | solicitar; procurar; buscar (cargo, salário, etc.) | beber até o fim | tratar com indiferença; ignorar}
  \seealsoref{干儿}{gan1'er2}
  \seealsoref{干儿}{gan1r5}
  \antonymref{潮}{chao2}
  \antonymref{湿}{shi1}
  \antonymref{鲜}{xian1}
  \end{Phonetics}
  \begin{Phonetics}{干}{gan4}[][HSK 1]
    \definition{adj.}{capaz; competente; habilidoso}
    \definition{s.}{tronco; parte principal; corpo principal ou parte importante de algo | habilidade; capacidade; competência}
    \definition{v.}{fazer; trabalhar; cuidar; fazer coisas | ocupar o cargo de; estar envolvido em; assumir, exercer | lutar; golpear; esforçar"-se}
  \synonymref{从}{cong2}
  \synonymref{搞}{gao3}
  \synonymref{任}{ren4}
  \synonymref{糟}{zao1}
  \synonymref{做}{zuo4}
  \antonymref{潮}{chao2}
  \antonymref{好}{hao3}
  \antonymref{湿}{shi1}
  \antonymref{鲜}{xian1}
  \antonymref{支}{zhi1}
  \end{Phonetics}
\end{Entry}

\begin{Entry}{干儿}{3,2}{⼲,⼉}
  \begin{Phonetics}{干儿}{gan1'er2}
    \definition{s.}{filho adotivo (adoção tradicional, ou seja, sem implicações legais)}
  \end{Phonetics}
  \begin{Phonetics}{干儿}{gan1r5}
    \definition{s.}{alimentos secos, desidratados}
  \end{Phonetics}
\end{Entry}

\begin{Entry}{干与}{3,3}{⼲,⼀}
  \begin{Phonetics}{干与}{gan1yu4}
    \variantof{干预}
  \end{Phonetics}
\end{Entry}

\begin{Entry}{干什么}{3,4,3}{⼲,⼈,⼃}
  \begin{Phonetics}{干什么}{gan4shen2me5}[][HSK 1]
    \definition{adv.}{o que fazer; o que ele está fazendo?; o que você está fazendo?; perguntar a razão ou o objetivo}
  \synonymref{为何}{wei4he2}
  \synonymref{为什么}{wei4shen2me5}
  \end{Phonetics}
\end{Entry}

\begin{Entry}{干戈}{3,4}{⼲,⼽}
  \begin{Phonetics}{干戈}{gan1ge1}[][HSK 7-9]
    \definition{s.}{armas de guerra; armas; guerra}[化干戈为玉帛。===Transforme guerra em paz.]
  \antonymref{玉帛}{yu4bo2}
  \end{Phonetics}
\end{Entry}

\begin{Entry}{干吗}{3,6}{⼲,⼝}
  \begin{Phonetics}{干吗}{gan4ma2}[][HSK 3]
    \definition{pron.}{por que?}
    \definition{v.}{o que fazer?}
  \end{Phonetics}
\end{Entry}

\begin{Entry}{干你屁事}{3,7,7,8}{⼲,⼈,⼫,⼅}
  \begin{Phonetics}{干你屁事}{gan1 ni3 pi4shi4}
    \definition{interj.}{``Foda"-se!''; ``O que isso te importa?''}
  \end{Phonetics}
\end{Entry}

\begin{Entry}{干扰}{3,7}{⼲,⼿}
  \begin{Phonetics}{干扰}{gan1rao3}[][HSK 5]
    \definition{v.}{perturbar; incomodar | interferir; interromper o funcionamento adequado de equipamentos eletrônicos com sinais eletrônicos dispersos}
  \seealsoref{打扰}{da3rao3}
  \synonymref{打岔}{da3//cha4}
  \synonymref{打搅}{da3jiao3}
  \synonymref{干预}{gan1yu4}
  \synonymref{扰乱}{rao3luan4}
  \synonymref{骚扰}{sao1rao3}
  \synonymref{阻挠}{zu3nao2}
  \antonymref{帮助}{bang1zhu4}
  \antonymref{辅助}{fu3zhu4}
  \antonymref{配合}{pei4he2}
  \end{Phonetics}
\end{Entry}

\begin{Entry}{干旱}{3,7}{⼲,⽇}
  \begin{Phonetics}{干旱}{gan1han4}[][HSK 7-9]
    \definition{adj.}{árido; seco}
  \end{Phonetics}
\end{Entry}

\begin{Entry}{干事}{3,8}{⼲,⼅}
  \begin{Phonetics}{干事}{gan4shi5}[][HSK 7-9]
    \definition{s.}{secretário (ou funcionário administrativo) encarregado de algo | administrador | secretária executiva}
  \end{Phonetics}
\end{Entry}

\begin{Entry}{干净}{3,8}{⼲,⼎}
  \begin{Phonetics}{干净}{gan1jing4}[][HSK 1]
    \definition{adj.}{limpo; limpo e arrumado; sem poeira, impurezas, etc.}
    \definition{adv.}{completamente; totalmente; sem deixar nada para trás}
  \seealsoref{纯洁}{chun2jie2}
  \seealsoref{清洁}{qing1jie2}
  \synonymref{洁净}{jie2jing4}
  \synonymref{清爽}{qing1shuang3}
  \synonymref{洗净}{xi3jing4}
  \synonymref{整洁}{zheng3jie2}
  \antonymref{灰尘}{hui1chen2}
  \antonymref{污秽}{wu1hui4}
  \antonymref{犹豫}{you2yu4}
  \end{Phonetics}
\end{Entry}

\begin{Entry}{干杯}{3,8}{⼲,⽊}
  \begin{Phonetics}{干杯}{gan1//bei1}[][HSK 2]
    \definition{interj.}{`Saúde!''}
    \definition{v.+compl.}{fazer um brinde;  brindar até a última gota}
  \synonymref{祝酒}{zhu4jiu3}
  \end{Phonetics}
\end{Entry}

\begin{Entry}{干活}{3,9}{⼲,⽔}
  \begin{Phonetics}{干活}{gan4//huo2}
    \definition{v.+compl.}{trabalhar; trabalhar em um emprego; fazer algo que exige esforço físico ou mental, especialmente trabalho árduo ou laborioso}
  \synonymref{工作}{gong1zuo4}
  \synonymref{劳动}{lao2dong5}
  \antonymref{歇息}{xie1xi5}
  \antonymref{休憩}{xiu1qi4}
  \antonymref{休息}{xiu1xi5}
  \end{Phonetics}
\end{Entry}

\begin{Entry}{干活儿}{3,9,2}{⼲,⽔,⼉}
  \begin{Phonetics}{干活儿}{gan4huo2r5}[][HSK 2]
    \definition{v.}{trabalhar; gastar energia física ou mental para fazer algo, especialmente trabalho árduo ou esforçado.}
  \end{Phonetics}
\end{Entry}

\begin{Entry}{干涉}{3,10}{⼲,⽔}
  \begin{Phonetics}{干涉}{gan1she4}[][HSK 6]
    \definition{s.}{interferência; refere"-se ao ato ou comportamento de interferir nos assuntos dos outros}
    \definition{v.}{interferir; intervir; intrometer"-se; pedir ou impedir algo geralmente significa interferir quando não se deve}
  \seealsoref{干预}{gan1yu4}
  \synonymref{插手}{cha1//shou3}
  \synonymref{关联}{guan1lian2}
  \synonymref{关系}{guan1xi5}
  \synonymref{牵涉}{qian1she4}
  \antonymref{放任}{fang4ren4}
  \antonymref{任凭}{ren4ping2}
  \antonymref{听凭}{ting1ping2}
  \end{Phonetics}
\end{Entry}

\begin{Entry}{干脆}{3,10}{⼲,⾁}
  \begin{Phonetics}{干脆}{gan1cui4}[][HSK 5]
    \definition{adj.}{claro; direto; (falar, fazer coisas) sem hesitação; atitude clara}
    \definition{adv.}{justamente; diretamente; sem maiores considerações | usado na linguagem falada}
  \seealsoref{索性}{suo3xing4}
  \synonymref{果断}{guo3duan4}
  \synonymref{简洁}{jian3jie2}
  \synonymref{精炼}{jing1lian4}
  \synonymref{爽快}{shuang3kuai5}
  \synonymref{痛快}{tong4kuai4}
  \synonymref{直接}{zhi2jie1}
  \antonymref{犹豫}{you2yu4}
  \end{Phonetics}
\end{Entry}

\begin{Entry}{干部}{3,10}{⼲,⾢}
  \begin{Phonetics}{干部}{gan4bu4}[][HSK 7-9]
    \definition[名,个,位]{s.}{quadro; funcionário público; funcionário do governo; servidores públicos (excluindo soldados e pessoal geral) em órgãos estatais, militares e organizações populares; refere"-se àqueles que desempenham determinado trabalho de liderança ou gestão | líder; pessoal que ocupa determinados cargos de liderança}
  \end{Phonetics}
\end{Entry}

\begin{Entry}{干预}{3,10}{⼲,⾴}
  \begin{Phonetics}{干预}{gan1yu4}[][HSK 5]
    \definition{s.}{intromissão; intervenção | tem significado neutro}
    \definition{v.}{intrometer"-se; intervir; interpor"-se}
  \seealsoref{干涉}{gan1she4}
  \synonymref{干扰}{gan1rao3}
  \antonymref{任凭}{ren4ping2}
  \antonymref{听凭}{ting1ping2}
  \end{Phonetics}
\end{Entry}

\begin{Entry}{干燥}{3,17}{⼲,⽕}
  \begin{Phonetics}{干燥}{gan1zao4}[][HSK 7-9]
    \definition{adj.}{seco; árido; sem umidade ou muito pouca umidade | enfadonho; desinteressante; chato, sem graça}
  \end{Phonetics}
\end{Entry}

%%%%%%%%%% 平 %%%%%%%%%%
\subsection*{平}\addcontentsline{loh}{figure}{平}

\begin{Entry}{平}{5}{⼲}
  \begin{Phonetics}{平}{ping2}[][HSK 2]
    \definition*{s.}{Sobrenome: Ping}
    \definition{adj.}{plano; nivelado; uniforme; liso | igual; justo | mesma pontuação; empatado | médio; comum | silencioso; tranquilo | no mesmo nível; altura igual; sem diferença | imparcial; médio; equitativo | calmo; estável; tranquilo | comum;  frequente}
    \definition{s.}{no mesmo nível; em pé de igualdade com; igual | tom nivelado, um dos quatro tons do chinês clássico}
    \definition{v.}{tornar nivelado ou uniforme; nivelar | reprimir; suprimir | acalmar; tornar pacífico; silenciar (acalmar); conter a raiva | estar no mesmo nível | acalmar; amenizar; controlar a raiva}
  \synonymref{均}{jun1}
  \synonymref{镇}{zhen4}
  \antonymref{陡}{dou3}
  \end{Phonetics}
\end{Entry}

\begin{Entry}{平凡}{5,3}{⼲,⼏}
  \begin{Phonetics}{平凡}{ping2fan2}[][HSK 6]
    \definition{adj.}{comum; ordinário; normal; não surpreendente}
  \seealsoref{平庸}{ping2yong1}
  \synonymref{广泛}{guang3fan4}
  \synonymref{平常}{ping2chang2}
  \synonymref{平淡}{ping2dan4}
  \synonymref{普通}{pu3tong1}
  \synonymref{通常}{tong1chang2}
  \synonymref{通俗}{tong1su2}
  \synonymref{庸俗}{yong1su2}
  \antonymref{出色}{chu1se4}
  \antonymref{典型}{dian3xing2}
  \antonymref{杰出}{jie2chu1}
  \antonymref{奇妙}{qi2miao4}
  \antonymref{特殊}{te4shu1}
  \antonymref{突出}{tu1chu1}
  \antonymref{伟大}{wei3da4}
  \antonymref{珍贵}{zhen1gui4}
  \antonymref{著名}{zhu4ming2}
  \antonymref{卓越}{zhuo2yue4}
  \end{Phonetics}
\end{Entry}

\begin{Entry}{平方}{5,4}{⼲,⽅}
  \begin{Phonetics}{平方}{ping2fang1}[][HSK 4]
    \definition{s.}{Matemática: segunda potência (de uma quantidade); quadrado | metro quadrado (m²)}[那间房有十二平方。===O quarto tem doze metros quadrados.]
  \synonymref{面积}{mian4ji1}
  \end{Phonetics}
\end{Entry}

\begin{Entry}{平方市丈}{5,4,5,3}{⼲,⽅,⼱,⼀}
  \begin{Phonetics}{平方市丈}{ping2fang1 shi4zhang4}
    \definition{clas.}{pés quadrados}
  \end{Phonetics}
\end{Entry}

\begin{Entry}{平方米}{5,4,6}{⼲,⽅,⽶}
  \begin{Phonetics}{平方米}{ping2fang1mi3}[][HSK 6]
    \definition{s.}{metro quadrado; a unidade legal de medida de área, 1 metro quadrado é igual a 10.000 centímetros quadrados}
  \end{Phonetics}
\end{Entry}

\begin{Entry}{平日}{5,4}{⼲,⽇}
  \begin{Phonetics}{平日}{ping2ri4}[][HSK 7-9]
    \definition{adv.}{normalmente; em tempos normais; em dias comuns; em circunstâncias normais}
    \definition{s.}{dias úteis; dias comuns; dias da semana}
  \end{Phonetics}
\end{Entry}

\begin{Entry}{平台}{5,5}{⼲,⼝}
  \begin{Phonetics}{平台}{ping2tai2}[][HSK 6]
    \definition[个]{s.}{casa com telhado plano rebocado | terraço | plataforma móvel; metaforicamente, refere"-se às áreas, oportunidades, ambientes, espaços, etc. que fornecem suporte e garantia para algo | plataforma; um sistema em um computador eletrônico que consiste em software e hardware básicos; tal sistema pode suportar a execução de programas aplicativos e softwares aplicativos podem ser desenvolvidos nesse sistema | plataforma; lugar; falando metaforicamente, o mesmo nível ou grau}
  \end{Phonetics}
\end{Entry}

\begin{Entry}{平民}{5,5}{⼲,⽒}
  \begin{Phonetics}{平民}{ping2min2}[][HSK 7-9]
    \definition[个]{s.}{civis; pessoas comuns; refere"-se geralmente a pessoas comuns (em oposição a nobres ou classes privilegiadas)}
  \end{Phonetics}
\end{Entry}

\begin{Entry}{平价}{5,6}{⼲,⼈}
  \begin{Phonetics}{平价}{ping2jia4}[][HSK 7-9]
    \definition{s.}{preços estabilizados (ou normalizados, moderados) | paridade; preço justo; preço de paridade | preço razoável; preço baixo; preço justo | preço justo (estatal); preço fixado pelo Estado; preço normal, racional; o padrão da moeda oficial de um país; a taxa de câmbio padrão entre as moedas do padrão"-ouro de dois países}
    \definition{v.}{estabilizar os preços}
  \end{Phonetics}
\end{Entry}

\begin{Entry}{平地}{5,6}{⼲,⼟}
  \begin{Phonetics}{平地}{ping2di4}
    \definition{s.}{terreno nivelado; terreno plano | planície; plano; gramado}
    \definition{v.}{nivelar o terreno; rastelar o solo até ficar liso | rastelar a terra; rastelar o solo}
  \synonymref{平原}{ping2yuan2}
  \antonymref{高山}{gao1shan1}
  \antonymref{山坡}{shan1po1}
  \antonymref{要么}{yao4me5}
  \end{Phonetics}
\end{Entry}

\begin{Entry}{平安}{5,6}{⼲,⼧}
  \begin{Phonetics}{平安}{ping2'an1}[][HSK 2]
    \definition{s.}{seguro; bem; sem contratempos; sem acidentes; são e salvo}
  \synonymref{安定}{an1ding4}
  \synonymref{安宁}{an1ning2}
  \synonymref{安全}{an1quan2}
  \synonymref{和平}{he2ping2}
  \synonymref{平和}{ping2he2}
  \synonymref{太平}{tai4ping2}
  \antonymref{危险}{wei1xian3}
  \end{Phonetics}
\end{Entry}

\begin{Entry}{平均}{5,7}{⼲,⼟}
  \begin{Phonetics}{平均}{ping2jun1}[][HSK 4]
    \definition{adj.}{igual; médio}
    \definition{s.}{média}
    \definition{v.}{calcular a média de um conjunto de números}
  \seealsoref{均匀}{jun1yun2}
  \synonymref{平衡}{ping2heng2}
  \end{Phonetics}
\end{Entry}

\begin{Entry}{平时}{5,7}{⼲,⽇}
  \begin{Phonetics}{平时}{ping2shi2}[][HSK 2]
    \definition{s.}{em tempos normais; em tempos comuns | em tempo de paz; refere"-se a períodos normais}
  \seealsoref{平日}{ping2ri4}
  \synonymref{平常}{ping2chang2}
  \synonymref{普通}{pu3tong1}
  \synonymref{日常}{ri4chang2}
  \synonymref{通常}{tong1chang2}
  \synonymref{往常}{wang3chang2}
  \end{Phonetics}
\end{Entry}

\begin{Entry}{平和}{5,8}{⼲,⼝}
  \begin{Phonetics}{平和}{ping2he2}[][HSK 7-9]
    \definition{adj.}{suave; ameno; moderado; plácido | (medicamento) leve}[这种药的药性平和。===Este medicamento tem um efeito leve.]
  \end{Phonetics}
\end{Entry}

\begin{Entry}{平坦}{5,8}{⼲,⼟}
  \begin{Phonetics}{平坦}{ping2tan3}[][HSK 5]
    \definition{adj.}{plano; uniforme; nivelado; liso; sem elevações ou depressões (referindo"-se principalmente ao relevo)}
  \antonymref{艰险}{jian1xian3}
  \antonymref{荆棘}{jing1ji2}
  \antonymref{起伏}{qi3fu2}
  \end{Phonetics}
\end{Entry}

\begin{Entry}{平面}{5,9}{⼲,⾯}
  \begin{Phonetics}{平面}{ping2mian4}[][HSK 7-9]
    \definition[个]{s.}{plano; superfície plana; duas dimensões}
  \end{Phonetics}
\end{Entry}

\begin{Entry}{平原}{5,10}{⼲,⼚}
  \begin{Phonetics}{平原}{ping2yuan2}[][HSK 5]
    \definition[片,个]{s.}{campo; planície; terreno plano e extenso}
  \synonymref{平地}{ping2di4}
  \antonymref{高原}{gao1yuan2}
  \antonymref{丘陵}{qiu1ling2}
  \end{Phonetics}
\end{Entry}

\begin{Entry}{平息}{5,10}{⼲,⼼}
  \begin{Phonetics}{平息}{ping2xi1}[][HSK 7-9]
    \definition{v.}{acalmar"-se; aquietar"-se; emoções negativas como conflito, disputa ou raiva diminuem ou cessam | suprimir; abafar; sufocar (uma rebelião, etc.); usar a força para resolver problemas}
  \end{Phonetics}
\end{Entry}

\begin{Entry}{平常}{5,11}{⼲,⼱}
  \begin{Phonetics}{平常}{ping2chang2}[][HSK 2]
    \definition{adj.}{comum; normal; ordinário; nada de especial}
    \definition{adv.}{normalmente; geralmente; como regra geral}
  \seealsoref{通常}{tong1chang2}
  \synonymref{平淡}{ping2dan4}
  \synonymref{平凡}{ping2fan2}
  \synonymref{平时}{ping2shi2}
  \synonymref{普通}{pu3tong1}
  \synonymref{日常}{ri4chang2}
  \synonymref{寻常}{xun2chang2}
  \synonymref{一般}{yi4ban1}
  \synonymref{正常}{zheng4chang2}
  \antonymref{独特}{du2te4}
  \antonymref{非常}{fei1chang2}
  \antonymref{可贵}{ke3gui4}
  \antonymref{奇怪}{qi2guai4}
  \antonymref{奇妙}{qi2miao4}
  \antonymref{神秘}{shen2mi4}
  \antonymref{突出}{tu1chu1}
  \antonymref{稀罕}{xi1han5}
  \antonymref{异常}{yi4chang2}
  \end{Phonetics}
\end{Entry}

\begin{Entry}{平常心}{5,11,4}{⼲,⼱,⼼}
  \begin{Phonetics}{平常心}{ping2chang2xin1}[][HSK 7-9]
    \definition[个]{s.}{mente calma; atitude calma e serena; compostura}
  \end{Phonetics}
\end{Entry}

\begin{Entry}{平庸}{5,11}{⼲,⼴}
  \begin{Phonetics}{平庸}{ping2yong1}
    \definition{adj.}{medíocre; indiferente; comum; normal, nada de especial}
  \seealsoref{平凡}{ping2fan2}
  \synonymref{平淡}{ping2dan4}
  \synonymref{无能}{wu2neng2}
  \synonymref{庸俗}{yong1su2}
  \antonymref{才华}{cai2hua2}
  \antonymref{出色}{chu1se4}
  \antonymref{非凡}{fei1fan2}
  \antonymref{高超}{gao1chao1}
  \antonymref{高尚}{gao1shang4}
  \antonymref{杰出}{jie2chu1}
  \antonymref{奇迹}{qi2ji4}
  \antonymref{伟大}{wei3da4}
  \antonymref{优秀}{you1xiu4}
  \antonymref{卓越}{zhuo2yue4}
  \end{Phonetics}
\end{Entry}

\begin{Entry}{平淡}{5,11}{⼲,⽔}
  \begin{Phonetics}{平淡}{ping2dan4}[][HSK 7-9]
    \definition{adj.}{sem graça; monótono; entediante; insípido; prosaico; banal; desinteressante; insosso e sem graça; ordinário; sem qualquer reviravolta}
  \end{Phonetics}
\end{Entry}

\begin{Entry}{平等}{5,12}{⼲,⽵}
  \begin{Phonetics}{平等}{ping2deng3}[][HSK 2]
    \definition{adj.}{igual; igualdade; refere"-se ao fato de as pessoas gozarem de tratamento igualitário nos aspectos sociais, políticos, econômicos e jurídicos}
  \seealsoref{平衡}{ping2heng2}
  \synonymref{同等}{tong2deng3}
  \synonymref{相等}{xiang1deng3}
  \antonymref{悬殊}{xuan2shu1}
  \end{Phonetics}
\end{Entry}

\begin{Entry}{平稳}{5,14}{⼲,⽲}
  \begin{Phonetics}{平稳}{ping2wen3}[][HSK 4]
    \definition{adj.}{firme; estável; suave e constante; sem oscilações ou flutuações}
  \synonymref{安定}{an1ding4}
  \synonymref{安稳}{an1wen3}
  \synonymref{稳定}{wen3ding4}
  \synonymref{稳固}{wen3gu4}
  \antonymref{波动}{bo1dong4}
  \antonymref{动荡}{dong4dang4}
  \antonymref{激烈}{ji1lie4}
  \antonymref{强烈}{qiang2lie4}
  \antonymref{危机}{wei1ji1}
  \antonymref{摇晃}{yao2huang4}
  \end{Phonetics}
\end{Entry}

\begin{Entry}{平静}{5,14}{⼲,⾭}
  \begin{Phonetics}{平静}{ping2jing4}[][HSK 4]
    \definition{adj.}{(humor, ambiente, etc.) calmo; quieto; pacífico; tranquilo}
  \seealsoref{安静}{an1jing4}
  \synonymref{安宁}{an1ning2}
  \synonymref{从容}{cong2rong2}
  \synonymref{寂静}{ji4jing4}
  \synonymref{冷静}{leng3jing4}
  \synonymref{宁静}{ning2jing4}
  \synonymref{平和}{ping2he2}
  \synonymref{清静}{qing1jing4}
  \synonymref{坦然}{tan3ran2}
  \synonymref{稳定}{wen3ding4}
  \antonymref{烦躁}{fan2zao4}
  \antonymref{沸腾}{fei4teng2}
  \antonymref{混乱}{hun4luan4}
  \antonymref{激动}{ji1dong4}
  \antonymref{激烈}{ji1lie4}
  \antonymref{焦急}{jiao1ji2}
  \antonymref{强烈}{qiang2lie4}
  \antonymref{热闹}{re4nao5}
  \antonymref{着急}{zhao2//ji2}
  \antonymref{震撼}{zhen4han4}
  \end{Phonetics}
\end{Entry}

\begin{Entry}{平衡}{5,16}{⼲,⾏}
  \begin{Phonetics}{平衡}{ping2heng2}[][HSK 6]
    \definition{adj.}{balanceado; equilibrado; os aspectos opostos são iguais ou compensados em quantidade ou qualidade | equilibrado; várias forças atuam sobre um objeto com magnitude igual e direções opostas para manter o objeto estável}
    \definition{v.}{equilibrar; trazer ou manter em equilíbrio; tornar as coisas ou alimentos iguais em quantidade, qualidade ou força}
  \seealsoref{平等}{ping2deng3}
  \synonymref{对称}{dui4chen4}
  \synonymref{均衡}{jun1heng2}
  \synonymref{均匀}{jun1yun2}
  \synonymref{平均}{ping2jun1}
  \synonymref{平稳}{ping2wen3}
  \antonymref{倾向}{qing1xiang4}
  \antonymref{倾斜}{qing1xie2}
  \end{Phonetics}
\end{Entry}

%%%%%%%%%% 年 %%%%%%%%%%
\subsection*{年}\addcontentsline{loh}{figure}{年}

\begin{Entry}{年}{6}{⼲}
  \begin{Phonetics}{年}{nian2}[][HSK 1]
    \definition*{s.}{Sobrenome: Nian}
    \definition{clas.}{ano; usado para calcular o número de anos}
    \definition{s.}{ano | idade | um período (época) da história | colheita anual | Ano Novo | artigos para o dia de Ano Novo | um período da vida de uma pessoa; fases da vida humana divididas por idade}
  \seealsoref{岁}{sui4}
  \end{Phonetics}
\end{Entry}

\begin{Entry}{年代}{6,5}{⼲,⼈}
  \begin{Phonetics}{年代}{nian2dai4}[][HSK 3]
    \definition[个]{s.}{idade; anos; tempo; um período de tempo com características distintas na história | uma década de um século; período de dez anos}
  \synonymref{时代}{shi2dai4}
  \synonymref{时期}{shi2qi1}
  \synonymref{岁月}{sui4yue4}
  \end{Phonetics}
\end{Entry}

\begin{Entry}{年级}{6,6}{⼲,⽷}
  \begin{Phonetics}{年级}{nian2ji2}[][HSK 2]
    \definition[个]{s.}{série; ano; níveis divididos de acordo com o tempo de estudo dos alunos na escola}
  \synonymref{班级}{ban1ji2}
  \synonymref{学年}{xue2nian2}
  \end{Phonetics}
\end{Entry}

\begin{Entry}{年纪}{6,6}{⼲,⽷}
  \begin{Phonetics}{年纪}{nian2ji4}[][HSK 3]
    \definition[把,个]{s.}{idade (de uma pessoa)}
  \synonymref{年齿}{nian2chi3}
  \synonymref{年龄}{nian2ling2}
  \synonymref{寿命}{shou4ming4}
  \synonymref{岁数}{sui4shu4}
  \end{Phonetics}
\end{Entry}

\begin{Entry}{年迈}{6,6}{⼲,⾡}
  \begin{Phonetics}{年迈}{nian2mai4}[][HSK 7-9]
    \definition{adj.}{velho; idoso}
  \end{Phonetics}
\end{Entry}

\begin{Entry}{年初}{6,7}{⼲,⾐}
  \begin{Phonetics}{年初}{nian2chu1}[][HSK 3]
    \definition{s.}{o começo do ano; os primeiros dias do ano}
  \antonymref{年底}{nian2di3}
  \antonymref{年终}{nian2zhong1}
  \end{Phonetics}
\end{Entry}

\begin{Entry}{年夜饭}{6,8,7}{⼲,⼣,⾷}
  \begin{Phonetics}{年夜饭}{nian2ye4fan4}[][HSK 7-9]
    \definition[顿,桌]{s.}{jantar de família na véspera de Ano Novo; um jantar especial realizado na véspera de Ano Novo (na noite de 31 de dezembro)}
  \end{Phonetics}
\end{Entry}

\begin{Entry}{年底}{6,8}{⼲,⼴}
  \begin{Phonetics}{年底}{nian2di3}[][HSK 3]
    \definition[个]{s.}{fim de ano; o fim do ano; geralmente os últimos dias de dezembro ou o fim do ano}
  \synonymref{年终}{nian2zhong1}
  \antonymref{年初}{nian2chu1}
  \end{Phonetics}
\end{Entry}

\begin{Entry}{年画}{6,8}{⼲,⽥}
  \begin{Phonetics}{年画}{nian2hua4}[][HSK 7-9]
    \definition{s.}{fotos de Ano Novo (ou Festival da Primavera); durante o Ano Novo Lunar, são afixadas imagens que retratam alegria e prosperidade}
  \seealsoref{年画儿}{nian2hua4r5}
  \end{Phonetics}
\end{Entry}

\begin{Entry}{年画儿}{6,8,2}{⼲,⽥,⼉}
  \begin{Phonetics}{年画儿}{nian2hua4r5}
    \definition{s.}{foto de Ano Novo (Festival da Primavera)}
  \end{Phonetics}
\end{Entry}

\begin{Entry}{年终}{6,8}{⼲,⽷}
  \begin{Phonetics}{年终}{nian2zhong1}[][HSK 7-9]
    \definition{s.}{fim de ano; fim do ano}
  \end{Phonetics}
\end{Entry}

\begin{Entry}{年货}{6,8}{⼲,⾙}
  \begin{Phonetics}{年货}{nian2huo4}
    \definition{s.}{artigos para o Ano Novo; suprimentos para o Festival da Primavera; itens tradicionalmente usados durante o Ano Novo Lunar, como doces, imagens de Ano Novo e fogos de artifício}
  \end{Phonetics}
\end{Entry}

\begin{Entry}{年限}{6,8}{⼲,⾩}
  \begin{Phonetics}{年限}{nian2xian4}[][HSK 7-9]
    \definition{s.}{limite de idade; número fixo de anos; o número de anos especificado ou usado como padrão geral}
  \end{Phonetics}
\end{Entry}

\begin{Entry}{年齿}{6,8}{⼲,⿒}
  \begin{Phonetics}{年齿}{nian2chi3}
    \definition{s.}{Literário: idade}[年齿尚幼。===Ainda jovem.]
  \synonymref{年纪}{nian2ji4}
  \synonymref{年龄}{nian2ling2}
  \synonymref{岁数}{sui4shu4}
  \end{Phonetics}
\end{Entry}

\begin{Entry}{年前}{6,9}{⼲,⼑}
  \begin{Phonetics}{年前}{nian2qian2}[][HSK 5]
    \definition{s.}{(pouco) antes da virada do ano | antes do final do ano | antes do ano novo}
  \end{Phonetics}
\end{Entry}

\begin{Entry}{年度}{6,9}{⼲,⼴}
  \begin{Phonetics}{年度}{nian2du4}[][HSK 5]
    \definition{s.}{ano; de acordo com a natureza e as necessidades de um negócio, há um prazo de doze meses com data de início e término definidas}
  \end{Phonetics}
\end{Entry}

\begin{Entry}{年轻}{6,9}{⼲,⾞}
  \begin{Phonetics}{年轻}{nian2qing1}[][HSK 2]
    \definition{adj.}{jovem; não muito velho (geralmente se refere a pessoas entre 10 e 20 anos)}
  \synonymref{青春}{qing1chun1}
  \antonymref{年迈}{nian2mai4}
  \antonymref{衰老}{shuai1lao3}
  \end{Phonetics}
\end{Entry}

\begin{Entry}{年龄}{6,13}{⼲,⿒}
  \begin{Phonetics}{年龄}{nian2ling2}[][HSK 5]
    \definition[个,段]{s.}{idade; animais, plantas e outros seres vivos vivem e crescem no mundo durante um determinado número de anos}
  \synonymref{年齿}{nian2chi3}
  \synonymref{年纪}{nian2ji4}
  \synonymref{寿命}{shou4ming4}
  \synonymref{岁数}{sui4shu4}
  \end{Phonetics}
\end{Entry}

\begin{Entry}{年薪}{6,16}{⼲,⾋}
  \begin{Phonetics}{年薪}{nian2xin1}[][HSK 7-9]
    \definition{s.}{salário anual; remuneração anual}[他的年薪已经达到了五十万元。===Seu salário anual chegou a 500.000 yuans.]
  \end{Phonetics}
\end{Entry}

%%%%%%%%%% 并 %%%%%%%%%%
\subsection*{并}\addcontentsline{loh}{figure}{并}

\begin{Entry}{并}{6}{⼲}
  \begin{Phonetics}{并}{bing1}
    \definition*{s.}{Bing, região originalmente parte de Shanxi (山西) | uma das nove províncias antigas da China (中国), geralmente referindo"-se às atuais Baoding (保定) e Zhengding (正定), na província de Hebei (河北), e Taiyuan (太原) e Datong (大同), na província de Shanxi (山西)}
  \seealsoref{保定}{bao3ding4}
  \seealsoref{大同}{da4tong2}
  \seealsoref{河北}{he2bei3}
  \seealsoref{山西}{shan1xi1}
  \seealsoref{太原}{tai4yuan2}
  \seealsoref{正定}{zheng4ding4}
  \seealsoref{中国}{zhong1guo2}
  \end{Phonetics}
  \begin{Phonetics}{并}{bing4}[][HSK 3,4]
    \definition{adv.}{lado a lado; igualmente; simultaneamente | (usado para reforçar uma negação) na verdade; definitivamente | mesmo assim | (usado para reforçar uma negação) na verdade; de forma alguma | todos; indica o conjunto completo, equivalente a 全部}
    \definition{conj.}{e; além disso}
    \definition{v.}{combinar; fundir; incorporar | ficar (ou colocar) lado a lado | estar paralelo a | anexar; juntar}
  \seealsoref{并且}{bing4qie3}
  \seealsoref{全部}{quan2bu4}
  \synonymref{合}{he2}
  \synonymref{且}{qie3}
  \end{Phonetics}
\end{Entry}

\begin{Entry}{并且}{6,5}{⼲,⼀}
  \begin{Phonetics}{并且}{bing4qie3}[][HSK 3]
    \definition{conj.}{e; bem como; usado entre verbos, adjetivos ou frases paralelas para indicar que várias ações são realizadas ao mesmo tempo ou que propriedades existem ao mesmo tempo | além disso; além do mais; ademais; usado na segunda metade de uma frase complexa para expressar um significado adicional}
  \seealsoref{并}{bing4}
  \seealsoref{而且}{er2qie3}
  \synonymref{同时}{tong2shi2}
  \synonymref{再者}{zai4zhe3}
  \end{Phonetics}
\end{Entry}

\begin{Entry}{并列}{6,6}{⼲,⼑}
  \begin{Phonetics}{并列}{bing4lie4}[][HSK 7-9]
    \definition{v.}{ficar lado a lado; ser justaposto; ter a mesma importância; organizar lado a lado}
  \end{Phonetics}
\end{Entry}

\begin{Entry}{并行}{6,6}{⼲,⾏}
  \begin{Phonetics}{并行}{bing4xing2}[][HSK 7-9]
    \definition{adj.}{simultâneo | Computação: paralelo | lado a lado (de dois processos, desenvolvimentos, pensamentos etc.)}
    \definition{v.}{caminhar lado a lado; correr lado a lado | fazer duas coisas ao mesmo tempo | prosseguir em paralelo}
  \end{Phonetics}
\end{Entry}

\begin{Entry}{并购}{6,8}{⼲,⾙}
  \begin{Phonetics}{并购}{bing4gou4}[][HSK 7-9]
    \definition{s.}{aquisição; fusão e aquisição}
    \definition{v.}{fundir; adquirir | assumir}
  \end{Phonetics}
\end{Entry}

\begin{Entry}{并非}{6,8}{⼲,⾮}
  \begin{Phonetics}{并非}{bing4fei1}[][HSK 7-9]
    \definition{adv.}{realmente não é; na verdade}
  \end{Phonetics}
\end{Entry}

\begin{Entry}{并排}{6,11}{⼲,⼿}
  \begin{Phonetics}{并排}{bing4pai2}
    \definition*{adv.}{lado a lado}
  \seealsoref{并列}{bing4lie4}
  \end{Phonetics}
\end{Entry}

%%%%%%%%%% 幷 %%%%%%%%%%
\subsection*{幷}\addcontentsline{loh}{figure}{幷}

\begin{Entry}{幷}{8}{⼲}
  \begin{Phonetics}{幷}{bing4}
    \variantof{并}
  \end{Phonetics}
\end{Entry}

%%%%%%%%%% 幸 %%%%%%%%%%
\subsection*{幸}\addcontentsline{loh}{figure}{幸}

\begin{Entry}{幸}{8}{⼲}
  \begin{Phonetics}{幸}{xing4}
    \definition*{s.}{Sobrenome: Xing}
    \definition{adj.}{feliz}
    \definition{adv.}{afortunadamente; felizmente}
    \definition{s.}{felicidade}
    \definition{v.}{alegrar"-se; sentir"-se feliz e contente | favorecer; patrocinar | vir; chegar; antigamente, referia"-se à chegada de um monarca a um determinado lugar}
  \end{Phonetics}
\end{Entry}

\begin{Entry}{幸亏}{8,3}{⼲,⼆}
  \begin{Phonetics}{幸亏}{xing4kui1}[][HSK 7-9]
    \definition{adv.}{felizmente; por sorte; por obra do destino; isso indica que consequências adversas foram evitadas por acaso devido a certas condições}
  \seealsoref{多亏}{duo1kui1}
  \seealsoref{幸好}{xing4hao3}
  \synonymref{好在}{hao3zai4}
  \end{Phonetics}
\end{Entry}

\begin{Entry}{幸好}{8,6}{⼲,⼥}
  \begin{Phonetics}{幸好}{xing4hao3}[][HSK 7-9]
    \definition{adv.}{felizmente; por sorte; afortunadamente}
  \synonymref{多亏}{duo1kui1}
  \synonymref{好在}{hao3zai4}
  \synonymref{幸亏}{xing4kui1}
  \antonymref{可惜}{ke3xi1}
  \end{Phonetics}
\end{Entry}

\begin{Entry}{幸存}{8,6}{⼲,⼦}
  \begin{Phonetics}{幸存}{xing4cun2}[][HSK 7-9]
    \definition{v.}{sobreviver}
  \synonymref{幸免}{xing4mian3}
  \antonymref{遇难}{yu4//nan4}
  \end{Phonetics}
\end{Entry}

\begin{Entry}{幸免}{8,7}{⼲,⼉}
  \begin{Phonetics}{幸免}{xing4mian3}[][HSK 7-9]
    \definition{v.}{escapar por pura sorte; ter um escape por pouco; evitar por acaso}
  \synonymref{幸存}{xing4cun2}
  \end{Phonetics}
\end{Entry}

\begin{Entry}{幸运}{8,7}{⼲,⾡}
  \begin{Phonetics}{幸运}{xing4yun4}[][HSK 3]
    \definition{adj.}{sortudo; feliz; afortunado}
    \definition[个,点,丝]{s.}{boa sorte; boa fortuna}
  \seealsoref{侥幸}{jiao3xing4}
  \seealsoref{荣幸}{rong2xing4}
  \synonymref{好运}{hao3yun4}
  \synonymref{庆幸}{qing4xing4}
  \synonymref{幸好}{xing4hao3}
  \synonymref{运气}{yun4qi5}
  \antonymref{不幸}{bu2xing4}
  \antonymref{倒霉}{dao3//mei2}
  \antonymref{厄运}{e4yun4}
  \antonymref{可怜}{ke3lian2}
  \end{Phonetics}
\end{Entry}

\begin{Entry}{幸运儿}{8,7,2}{⼲,⾡,⼉}
  \begin{Phonetics}{幸运儿}{xing4yun4'er2}
    \definition{s.}{o favorito da fortuna; sujeito sortudo | cão sortudo; pessoa sortuda}
  \end{Phonetics}
\end{Entry}

\begin{Entry}{幸运抽奖}{8,7,8,9}{⼲,⾡,⼿,⼤}
  \begin{Phonetics}{幸运抽奖}{xing4yun4 chou1jiang3}
    \definition{s.}{loteria | sorteio}
  \end{Phonetics}
\end{Entry}

\begin{Entry}{幸福}{8,13}{⼲,⽰}
  \begin{Phonetics}{幸福}{xing4fu2}[][HSK 3]
    \definition{adj.}{feliz; a vida, a família e outras circunstâncias deixam as pessoas satisfeitas e felizes}
    \definition{s.}{felicidade; bem estar; sensação ou experiência satisfatória e feliz, etc.}
  \seealsoref{福}{fu2}
  \synonymref{快乐}{kuai4le4}
  \synonymref{美满}{mei3man3}
  \synonymref{甜蜜}{tian2mi4}
  \antonymref{悲哀}{bei1'ai1}
  \antonymref{悲惨}{bei1can3}
  \antonymref{悲伤}{bei1shang1}
  \antonymref{不幸}{bu2xing4}
  \antonymref{可怜}{ke3lian2}
  \antonymref{苦难}{ku3nan4}
  \antonymref{凄惨}{qi1can3}
  \antonymref{痛苦}{tong4ku3}
  \antonymref{灾难}{zai1nan4}
  \end{Phonetics}
\end{Entry}

%%%%% EOF %%%%%

